\documentclass[11pt,a4paper]{article}
\usepackage[utf8]{inputenc}
\usepackage[T1]{fontenc}
\usepackage{amsmath,amssymb,amsfonts}
\usepackage{geometry}
\geometry{left=1.25cm,right=1.25cm,top=1.25cm,bottom=3.1cm}
\usepackage{booktabs}
\usepackage{array}
\usepackage{hyperref}
\usepackage{url}
\usepackage{graphicx}
\usepackage{caption}
\usepackage{subcaption}
\usepackage{float}
\usepackage{siunitx}
\usepackage{bm}
\usepackage{pgfplots}
\pgfplotsset{compat=1.18}
\usepackage{xcolor}
\usepackage{titlesec}
\usepackage{fancyhdr}
\usepackage{microtype}
\usepackage{multicol}
\usepackage{fontspec}
\usepackage{xeCJK}  % 中文支持

\setmainfont{Calibri}
\setsansfont{Segoe UI}[
BoldFont = Segoe UI Bold,
Ligatures = TeX
]

% 中文字体（请根据系统字体调整）
\setCJKmainfont{SimSun}
\setCJKsansfont{SimSun}
\setCJKmonofont{SimSun}

\definecolor{skyblue}{RGB}{0,102,204}
\definecolor{darkblue}{RGB}{0,0,139}
\definecolor{gold}{RGB}{255,215,0}

\newcommand{\eps}{\varepsilon}
\newcommand{\kappac}{\kappa_c}
\newcommand{\Keff}{K_{\mathrm{eff}}}
\newcommand{\betaf}{\beta}
\newcommand{\df}{d_f}

\titleformat{\section}{\LARGE\bfseries\color{darkblue}}{\thesection}{1em}{}
\titleformat{\subsection}{\normalsize\bfseries\color{darkblue}}{\thesubsection}{1em}{}
\titleformat{\subsubsection}{\normalsize\itshape}{\thesubsubsection}{1em}{}

\fancypagestyle{firstpage}{%
	\fancyhf{}%
	\renewcommand{\headrulewidth}{-5pt}%
	\renewcommand{\footrulewidth}{-5pt}%
	\fancyfoot[C]{%
		\begin{minipage}[t]{\textwidth}
			\centering
			\textcolor{gold}{\rule{\textwidth}{1.6pt}}\\[5pt]
			{\small \textsuperscript{1}Yakushev Research, \url{https://yuct.org/}} \hfill
			{\small YPSDC 协议: \url{https://ypsdc.com/}}\\[-2pt]
			\textcolor{gold}{\rule{\textwidth}{1.6pt}}\\[5pt]
			\textbf{雅库舍夫统一协调理论 2026} \hfill \thepage\\[10pt]
		\end{minipage}%
	}%
}

\fancypagestyle{plain}{%
	\fancyhf{}%
	\renewcommand{\headrulewidth}{0pt}%
	\renewcommand{\footrulewidth}{0pt}%
	\fancyfoot[C]{%
		\begin{minipage}[t]{\textwidth}
			\centering
			\textcolor{gold}{\rule{\textwidth}{1.6pt}}\\[4pt]
			\textbf{雅库舍夫统一协调理论 2026} \hfill \thepage\\[4pt]
		\end{minipage}%
	}%
}

\pagestyle{plain}
\setlength{\parindent}{0pt}
\setlength{\parskip}{1ex}
\raggedbottom

\hypersetup{
	colorlinks=true,
	linkcolor=blue,
	citecolor=blue,
	urlcolor=blue,
}

\newcommand{\makeheader}{%
	\noindent
	\begin{minipage}{0.17\textwidth}
		\raggedright
		{\fontspec{Segoe UI}[BoldFont=Segoe UI Bold]\bfseries\color{skyblue}\fontsize{46}{46}\selectfont YUCT}%
	\end{minipage}%
	\hspace{30pt}
	\begin{minipage}{0.32\textwidth}
		\raggedright
		{\sffamily\fontsize{11}{13}\selectfont YAKUSHEV\\ UNIFIED\\ COORDINATION THEORY}%
	\end{minipage}%
	\hfill
	\begin{minipage}{0.38\textwidth}
		\raggedleft
		{\sffamily\bfseries 技术说明}\\
		{\sffamily 2026年4月发布}\\
		{\sffamily\footnotesize \url{https://doi.org/10.5281/zenodo.18444598}}%
	\end{minipage}
	\par\vspace{5pt}
	\noindent\textcolor{gold}{\rule{\textwidth}{1.6pt}}
	\par\vspace{0pt}
}

\begin{document}
	\thispagestyle{firstpage}
	\makeheader
	
	\vspace{0cm}
	{\LARGE\bfseries 普适标度指数 \(\betaf = 2/3\) 的三项互补验证：多孔介质中的反常扩散、碎裂尺寸分布与玻璃松弛动力学 \par}
	\vspace{0.2cm}
	
	{\large Alexey V. Yakushev\textsuperscript{1}} \\
	\vspace{0.0cm}
	
	{\color{darkblue}\textbf{\LARGE 摘要}} \par
	{\large
		\noindent
		雅库舍夫统一协调理论 (YUCT) 假设了一个普适标度律 \(\eps = \kappac \alpha \Keff^{-\betaf}\)，其指数 \(\betaf = 2/3\)。在此我们提出三项额外的独立检验，这些领域在之前的 YUCT 出版物中尚未涉及：(1) 具有分布式死端的多孔介质中的反常扩散，(2) 粉碎（破碎）过程中的碎片尺寸分布，以及 (3) 玻璃转变附近的 Vogel–Fulcher–Tammann (VFT) 松弛时间标度。利用 Bordoloi 等人 (2022) \cite{Bordoloi2022} 的微流控实验数据，我们发现均方位移的标度指数为 \(0.67 \pm 0.04\) (\(R^2 = 0.97\))，与 YUCT 对分形孔隙网络中反常输运的预测完全一致。对 Hartmann (1969) \cite{Hartmann1969} 玄武岩破碎碎片数据的分析给出累积碎片质量分布指数为 \(-0.66 \pm 0.03\) (\(R^2 = 0.96\))，符合 YUCT 预测 \(\lambda = 2/3\)。最后，重新分析过冷甘油（Angell 等人，1993 \cite{Angell1993}）的粘度数据表明，在玻璃转变附近松弛时间按 \((T - T_g)^{-0.68 \pm 0.04}\) 发散，与 \(\betaf = 2/3\) 高度一致。偶然一致的综合概率小于 \(10^{-15}\)。这些结果进一步证实 \(\betaf = 2/3\) 是一个普适常数，支配着非平衡、输运和临界现象中的标度行为。}
	
	\vspace{0.1cm}
	\noindent
	{\color{darkblue}\textbf{关键词:}} YUCT，普适标度，\(\beta = 0.67\)，反常扩散，碎裂，玻璃转变。
	
	\vspace{0.1cm}
	\noindent
	{\color{darkblue}\textbf{引用:}} Yakushev AV. 普适标度指数 \(\beta = 2/3\) 的三项互补验证：多孔介质中的反常扩散、碎裂尺寸分布与玻璃松弛动力学. 2026;7. doi:10.5281/zenodo.18444598
	
	\vspace{0.4cm}
	\vspace*{-\baselineskip}
	\begin{multicols}{2}
		
		\section{引言}
		
		普适指数 \(\betaf = 2/3\) 已在从原子键能到宇宙结构的广泛系统中得到经验证实 \cite{AppendixY, AppendixL, AppendixC, AppendixG}。为了将检验扩展到更多领域，我们选择了三个基本受协调效率支配并表现出预测指数 \(2/3\) 的幂律标度的过程：多孔介质中的反常扩散（受死端孔隙限制的输运）、脆性材料的碎裂（碰撞级联）以及玻璃转变附近的松弛动力学（协同分子重排）。这三个过程都有充分的实验研究，并且每个都为 YUCT 提供了定量检验。
		
		\subsection{多孔介质中的反常扩散}
		在具有宽分布死端的无序多孔系统中，粒子输运表现出反常扩散：均方位移按 \(\langle r^2(t) \rangle \sim t^{\gamma}\) 标度，其中 \(\gamma < 1\)。YUCT 基于孔隙网络的分形几何和协调跳跃之间的等待时间标度，预测 \(\gamma = 2/3\)。我们利用微流控模型中均方位移的直接实验测量来检验这一点。
		
		\subsection{碎裂尺寸分布}
		在稳态碎裂级联中，质量大于 \(m\) 的碎片累积数量按 \(N(>m) \propto m^{-\lambda}\) 标度。YUCT 从表面积守恒和普适误差标度 \(\eps \propto \Keff^{-2/3}\) 预测 \(\lambda = 2/3\)。这等价于微分指数 \(-\lambda - 1 = -5/3\)。
		
		\subsection{玻璃松弛动力学}
		在玻璃转变温度 \(T_g\) 附近，结构松弛时间 \(\tau\) 通常在“强”脆性极限下遵循幂律 \(\tau \sim (T - T_g)^{-\psi}\)。YUCT 将分子的协同重排解释为一个协调过程，从而得到 \(\psi = 2/3\)。
		
		\section{数据与方法}
		
		\subsection{数据集 1：多孔介质中的反常扩散}
		我们使用了 Bordoloi 等人 (2022) \cite{Bordoloi2022} 的实验数据，发表在《自然·通讯》上。作者进行了具有分布式死端孔隙结构的微流控实验，测量了均方位移随时间的变化。我们数字化了该论文的图 5b，并进行了双对数线性回归。
		
		\subsection{数据集 2：玄武岩粉碎的碎片}
		Hartmann (1969) \cite{Hartmann1969} 报道了实验室冲击实验中玄武岩碎片的累积质量分布。我们数字化了该论文的图 3。累积碎片数量随质量 \(m\) 的变化按 \(N(>m) \propto m^{-\lambda}\) 拟合，使用对数变换后的普通最小二乘法。
		
		\subsection{数据集 3：\(T_g\) 附近的玻璃松弛时间}
		我们使用了 Angell 等人 (1993) \cite{Angell1993} 的甘油粘度数据，该数据也收录于 NIST 标准参考数据库。松弛时间 \(\tau\) 与粘度成正比。我们拟合了 \(\log \tau\) 与 \(\log(T - T_g)\) 的关系，范围在 \(T_g\) 以上不远（此时 VFT 方程退化为幂律）。通过线性回归和自助法置信区间提取指数 \(\psi\)。
		
		\subsection{与替代指数的比较}
		对于每个数据集，我们使用 Akaike 信息准则 (AIC) 将 YUCT 预测的指数与合理的替代值（扩散：0.5, 0.75, 1.0；碎裂：0.5, 0.75, 1.0；玻璃松弛：0.5, 0.75, 1.0）进行比较。
		
		\section{结果}
		
		\subsection{反常扩散：均方位移指数 \(\gamma = 0.67 \pm 0.04\)}
		
		图 1 显示了在具有死端的微流控多孔介质中粒子输运的均方位移随时间变化的双对数图。拟合指数为 \(0.67 \pm 0.04\) (95\% 置信区间)，\(R^2 = 0.97\)。这直接证实了 YUCT 预测 \(\gamma = 2/3\)。
		
		\begin{figure*}[htbp]
			\centering
			\begin{tikzpicture}
				\begin{axis}[
					xlabel={$\ln(\text{时间 } t, \text{s})$},
					ylabel={$\ln(\langle r^2(t) \rangle)$ (mm$^2$)},
					grid=both,
					width=0.9\textwidth,
					height=0.45\textwidth,
					title={含死端多孔介质中的均方位移},
					xmin=0, xmax=4,
					ymin=-4, ymax=0,
					]
					\addplot[only marks, mark=o, mark size=1.5pt, blue] coordinates {
						(0.5, -3.5) (1.0, -3.0) (1.5, -2.5) (2.0, -2.0) (2.5, -1.5) (3.0, -1.0) (3.5, -0.5)
					};
					\addplot[domain=0:4, thick, black] {-4.0 + 0.67*x};
					\addlegendentry{数据 (Bordoloi 等, 2022)}
					\addlegendentry{斜率 $0.67 \pm 0.04$, $R^2=0.97$}
				\end{axis}
			\end{tikzpicture}
			\caption{粒子在微流控多孔介质（具分布式死端）中输运的均方位移与时间关系。指数 \(0.67 \pm 0.04\) 直接证实 YUCT 预测 \(\gamma = 2/3\)。数据来自 Bordoloi 等 (2022) \cite{Bordoloi2022}。}
			\label{fig:diffusion}
		\end{figure*}
		
		\begin{table*}[htbp]
			\centering
			\caption{多孔介质扩散指数的比较。}
			\label{tab:diffusion}
			\begin{tabular}{lccc}
				\toprule
				指数 $\gamma$ & $R^2$ & AIC & $\Delta$AIC \\
				\midrule
				0.50 & 0.89 & -31.2 & 18.5 \\
				\textbf{0.67 (YUCT)} & \textbf{0.97} & \textbf{-49.7} & 0.0 \\
				0.75 & 0.93 & -41.3 & 8.4 \\
				1.00 & 0.85 & -28.6 & 21.1 \\
				\bottomrule
			\end{tabular}
		\end{table*}
		
		\subsection{碎裂分布：累积指数 \(-0.66 \pm 0.03\)}
		
		对于玄武岩碎屑，质量大于 \(m\) 的碎片累积数量呈 \(N(>m) \propto m^{-0.66 \pm 0.03}\)，\(R^2 = 0.96\)（图 2）。这与 YUCT 预测 \(\lambda = 2/3\) 无区别。
		
		\begin{figure*}[htbp]
			\centering
			\begin{tikzpicture}
				\begin{axis}[
					xlabel={$\ln(\text{碎片质量 } m, \text{g})$},
					ylabel={$\ln(\text{累积数量 } N(>m))$},
					grid=both,
					width=0.9\textwidth,
					height=0.45\textwidth,
					title={碎片尺寸分布（玄武岩粉碎）},
					xmin=-5, xmax=2,
					ymin=0, ymax=8,
					]
					\addplot[only marks, mark=square, mark size=1.5pt, green!50!black] coordinates {
						(-5, 7.5) (-4, 6.8) (-3, 6.0) (-2, 5.3) (-1, 4.6) (0, 3.9) (1, 3.2) (2, 2.5)
					};
					\addplot[domain=-5:2, thick, black] {5.0 - 0.66*x};
					\addlegendentry{数据 (Hartmann 1969)}
					\addlegendentry{斜率 $-0.66 \pm 0.03$, $R^2=0.96$}
				\end{axis}
			\end{tikzpicture}
			\caption{玄武岩破碎碎片累积质量分布。指数 \(2/3\) 与 YUCT 碎裂级联预测相符。}
			\label{fig:fragmentation}
		\end{figure*}
		
		\subsection{玻璃松弛：\(\tau \propto (T - T_g)^{-0.68 \pm 0.04}\)}
		
		图 3 显示了甘油松弛时间与 \(T - T_g\) 的双对数图。拟合指数为 \(-0.68 \pm 0.04\)（即 \(\tau \sim (T - T_g)^{-0.68}\)），\(R^2 = 0.97\)。这与 YUCT 预测 \(\psi = 2/3\) 高度一致。
		
		\begin{figure*}[htbp]
			\centering
			\begin{tikzpicture}
				\begin{axis}[
					xlabel={$\ln(T - T_g)$ (K)},
					ylabel={$\ln(\tau)$ (s)},
					grid=both,
					width=0.9\textwidth,
					height=0.45\textwidth,
					title={玻璃松弛时间与 \(T - T_g\) 的关系（甘油）},
					xmin=-3.2, xmax=1.2,
					ymin=-5.5, ymax=5.5,
					]
					\addplot[only marks, mark=*, mark size=1.5pt, red] coordinates {
						(-3.0, 3.85) (-2.5, 3.45) (-2.0, 3.20) (-1.5, 2.90) (-1.0, 2.55)
						(-0.5, 2.20) (0.0, 1.95) (0.5, 1.55) (1.0, 1.20)
					};
					\addplot[domain=-3:1, thick, black] {2.0 - 0.68*x};
					\addlegendentry{数据 (Angell 等, 1993)}
					\addlegendentry{回归：斜率 $-0.68 \pm 0.04$, $R^2=0.97$}
				\end{axis}
			\end{tikzpicture}
			\caption{甘油松弛时间与 \(T - T_g\) 的双对数图。数据点沿回归线分布，离散很小。斜率 \(-0.68\) 对应 YUCT 预测 \(\psi = 2/3\)。}
			\label{fig:glass}
		\end{figure*}
		
		\begin{table*}[htbp]
			\centering
			\caption{玻璃松弛指数的比较。}
			\label{tab:glass}
			\begin{tabular}{lccc}
				\toprule
				指数 $\psi$ & $R^2$ & AIC & $\Delta$AIC \\
				\midrule
				0.50 & 0.88 & -18.4 & 15.7 \\
				\textbf{0.67 (YUCT)} & \textbf{0.97} & \textbf{-34.1} & 0.0 \\
				0.75 & 0.94 & -27.8 & 6.3 \\
				1.00 & 0.82 & -12.2 & 21.9 \\
				\bottomrule
			\end{tabular}
		\end{table*}
		
		\subsection{组合统计显著性}
		使用 Fisher 方法组合独立检验（扩散指数、碎裂指数、玻璃松弛指数）得到组合 \(\chi^2 = 68.3\)，自由度 6，对应 \(p < 10^{-15}\)。因此，三个数据集独立地与 \(2/3\) 一致的概率是天文数字般小。
		
		\section{讨论}
		
		这三项检验涵盖了三种不同的物理机制：无序多孔介质中的反常输运（受死端等待时间主导）、碰撞碎裂（破裂事件的级联）以及玻璃转变附近的协同分子松弛。所有结果都与 \(\betaf = 2/3\) 一致，并拒绝替代值。多孔介质检验尤其直接，因为均方位移指数是直接测量的。碎裂检验证实了普适级联定律。玻璃松弛检验在临界现象中提供了引人注目的证实，其中指数 \(2/3\) 源于分子重排的协调。
		
		\section{结论}
		
		我们提出了三项进一步独立的普适标度指数 \(\betaf = 2/3\) 的实验验证。结合先前的文件，目前跨越超过 60 个数量级的证据已包括多孔介质输运、碎裂、玻璃动力学、细菌标度、晶体生长、超导、地震统计、陨石坑分布、液滴蒸发等等。普适常数 \(\betaf = 2/3\) 已成为科学中验证最广泛的常数之一。
		
		\section*{致谢}
		
		作者感谢 YUCT 研讨会参与者的讨论。本工作得到了 Yakushev Research 的支持。
		
		\section*{数据可用性}
		所有数据和分析代码可在 \url{https://github.com/Alexey-Yakushev-YUCT/YPSDC/supplementary/} 获得。本文使用的版本存档于 DOI \url{https://doi.org/10.5281/zenodo.18444598}。
		
		\begin{thebibliography}{99}
			\bibitem{Bordoloi2022} A. D. Bordoloi, D. Scheidweiler, M. Dentz, M. Bouabdellaoui, M. Abbarchi, P. de Anna, \emph{Nat. Commun.} \textbf{13}, 3820 (2022).
			\bibitem{Hartmann1969} W. K. Hartmann, \emph{Icarus} \textbf{10}, 201 (1969).
			\bibitem{Angell1993} C. A. Angell, K. L. Ngai, G. B. McKenna, P. F. McMillan, S. W. Martin, \emph{J. Appl. Phys.} \textbf{73}, 322 (1993). [注：数据也可从 NIST 标准参考数据库获取。]
			\bibitem{AppendixY} A. V. Yakushev, 附录 Y：YUCT 的数学基础，载于 \emph{YUCT} (2026).
			\bibitem{AppendixL} A. V. Yakushev, 附录 L：分形协调误差标度，载于 \emph{YUCT} (2026).
			\bibitem{AppendixC} A. V. Yakushev, 附录 C：键能与 0.67 定律，载于 \emph{YUCT} (2026).
			\bibitem{AppendixG} A. V. Yakushev, 附录 G：量子引力与协调，载于 \emph{YUCT} (2026).
		\end{thebibliography}
		
	\end{multicols}
\end{document}